\documentclass[10pt]{beamer}

%\usetheme{AnnArbor}
\usepackage[dvipsnames]{xcolor}
\usetheme{Bergen}
\usecolortheme[named=purple]{structure}
%\usecolortheme{monarca}
\setbeamertemplate{navigation symbols}{}%\insertframenumber}
\def\mcolor {ForestGreen}
%\addtobeamertemplate{title page}{

%}
% Personnalisation de la barre latérale gauche
%\setbeamersize{sidebar width left=20pt}
%\setbeamersize{sidebar width left=default}
%\setbeamertemplate{sidebar canvas left}[horizontal shading]
%[left=purple,right=white]

%\setbeamertemplate{sidebar canvas left}{TEST}
%\setbeamertemplate{headline}{
%	TEST
%  \leavevmode%
%  \hbox{%
%    \begin{beamercolorbox}[wd=\paperwidth,ht=2.5ex,dp=1ex,leftskip=1em]{section in head/foot}
%      \insertsection
%    \end{beamercolorbox}%
%  }
%}

%\setbeamersize{text margin left=0.5cm,text margin right=0.5cm}

\usepackage[T1]{fontenc}
\usepackage[utf8]{inputenc}
%\usepackage{pgfplots}
%\usepackage{tikz}
\usepackage{tikz-cd}
%\usepackage{xcolor}
%\usepgfplotslibrary{fillbetween}
%\pgfplotsset{compat=1.18}
%\usepackage{graphicx}

%\usepackage [
%n, % or lambda
%advantage ,
%operators ,
%sets ,
%adversary ,
%landau ,
%probability ,
%notions ,
%logic ,
%ff ,
%mm,
%primitives ,
%events ,
%complexity ,
%oracles ,
%asymptotics ,
%keys
%]{cryptocode}
\theoremstyle{plain}
\newtheorem{proposition}[theorem]{\translate{Proposition}}
\theoremstyle{example}
\newtheorem{remark}[theorem]{\translate{Remark}}
%\newtheorem{problem}[problem]{\translate{Problem}}

\author{Mathieu \textsc{Bouget}}
\title[Studying FESTA]{Studying encryption from supersingular torsion attacks}
\institute{UMPA - Lyon,\\Benjamin Wesolowski,\\ Maria Corte-Real Santos}
\date{July 10th}

\begin{document}
\frame[plain]{\titlepage}

%TODO pourquoi SS ?
%TODO show number slide
\setcounter{framenumber}{0}
\begin{frame}{$\insertframenumber -$Topic}

	\begin{itemize}
		\item Isogeny-based cryptography
		\item 
	Public encryption scheme FESTA proposed by Basso, Maino and Pope
	in \cite{FESTA}	
\item Security aspect
	\end{itemize}
\end{frame}
%\frame[plain]{\tableofcontents}
\section{Preliminaries}
\begin{frame}{$\insertframenumber -$Elliptic curves}
	\hfill\raisebox{-110pt}[0pt][0pt]{\makebox[0pt][r]{
	%GRAPHIQUE
	\includegraphics[width=4cm]{media/addition-2.png}
}}
\begin{itemize}
	\item Let $a,b\in\mathbb F_{p^2}$ with $p \neq 2,3$
	\item $4a^3+27b^2\neq 0$
	\item $E:y^2 = x^3+ax+b$
\end{itemize}
\vspace{1cm}
\begin{itemize}
	\item $(P,Q)\in E^2\longmapsto P+Q\in E\quad$ (fast)
	%\item $P+Q+R=0\iff  (P,Q,R) \text{ aligned}$
	\item $(E,+,0)$ commutative group
	%\item Computing $P+Q$ is fast
\end{itemize}
\end{frame}

\begin{frame}{$\insertframenumber -$Isogenies}
\begin{itemize}
	\item An isogeny is an algebraic morphism, ie \[
			\varphi :(x,y)\in E_1\mapsto
			\left(\frac{f_1(x,y)}{f_3(x,y)},\frac{f_2(x,y)}{f_3(x,y)}\right)\in E_2,
			f_1,f_2,f_3\in\mathbb F[X,Y]
	\] 
	%with $f_1,f_2,f_3$ polynomials in $\mathbb F[X,Y]$
\item with finite kernel $\ker\varphi$
\end{itemize}
\center\includegraphics[width=4cm]{media/dual-2.png}
\begin{itemize}
	\item $\varphi$ is a group morphism
	\item isogenies are stable under composition
	\item $\deg\varphi$ and $p$ coprime $\iff$  $\varphi$ separable $\iff$ $\deg\varphi =
		\#\ker\varphi$
	\item $\deg\varphi\circ\psi=\deg\varphi\deg\psi$
	\item $\varphi:P\in E\mapsto \varphi(P)\in E'$ takes $O(\deg\varphi)$ operations in
		general
	\item Any isogeny $\phi:E\to E'$ admits a dual isogeny $\widehat\phi:E'\to E$
	\item $\deg\widehat\phi=\deg\phi$
%	\item $\deg\varphi$ is the size of the kernel when $\varphi$ is \textbf{separable}
%	\item Given $\varphi:E_0\to E_1$ and $\psi:E_1\to E_2$ \begin{itemize}
%			\item $\psi\circ\varphi:E_0\to E_2$ is an isogeny
%			\item $\deg(\psi\circ\varphi)=\deg\varphi\deg\psi$
%		\end{itemize}
\end{itemize}
\end{frame}

\begin{frame}{$\insertframenumber -$$l$-isogeny graph}


	Let $l$ be a small prime, and $\varphi_1,\dots,\varphi_n$ be $l$-isogenies.\\
	%TODO ajouter figures avec voisins ?
	%TODO figure du graphe

	\begin{itemize}
		\item Vertices : Elliptic curves
		\item Edges : $l$-isogenies
		\item $(l+1)$-regular graph
	\end{itemize}
	\begin{center}\includegraphics[width=7cm]{media/ln-iso-2.png}\end{center}
	The isogeny $\varphi = \varphi_n\circ\dots\circ\varphi_1$ has degree $l^n$ but can be
	computed in $O(nl)$ instead of $O(l^n)$.
\uncover<2>{
	\begin{problem}
	Given $E_0,E_n$ and $l$, finding $\varphi$ should be difficult ! 
	\end{problem}
}
	
\end{frame}
\begin{frame}{$\insertframenumber -$SIDH}
	%\begin{itemize}
	%	\item Public starting curve $E$
	%	\item Alice generates $\varphi_A:E\to E_A$ and sends $E_A$ 
	%	\item Bob generates $\varphi_B:E\to E_B$ and sends $E_B$
	%	\item Alice computes $E_{AB}="\varphi_A(E_B)"$ 
	%	\item Bob computes $E_{BA}="\varphi_B(E_A)"$ 
	%\end{itemize}
	\begin{itemize}
		\item Public information : $E,E_A,E_B$
		\item Alice secret : $\varphi_A$
		\item Bob secret : $\varphi_B$
		\item Alice and Bob combined secret $E_{AB} = "\varphi_A"(E_B) = "\varphi_B"(E_A)$
	\end{itemize}
	\begin{center}\includegraphics[width=5cm]{media/SIDH-2.png}\end{center}
	\uncover<2>{
		\alert{
	%We need Bob to know $\varphi_A(\ker\varphi_B)$ and Alice to know
	%$\varphi_B(\ker\varphi_A)$
			Need torsion-point information :(
}
}
\end{frame}

%TODO bouger duale vers une slide avant en même temps qu'on parle du graphe de l iso

\begin{frame}{$\insertframenumber -$$N$-torsion}
		\begin{itemize}
			\item $[N]_E:S\in E\longmapsto S+\dots+S\in E$ is an isogeny 
			\item $E[N]=\ker([N]_E)$ is the $N$-torsion of $E$
			\item when $N$ and $p$ are coprime, $E[N]\cong \mathbb Z_N\times\mathbb Z_N$
		\end{itemize}	
		\vspace{1cm}
		If $P,Q$ is a basis of $E[N]$ and $S=[u]P + [v]Q$
		\[
			\varphi(S)=\varphi([u]P+[v]Q)=[u]\varphi(P)+[v]\varphi(Q)
		\] 
	%\begin{center}\includegraphics[width=5cm]{media/dual-2.png}\end{center}
	% graphique ajouter boucle avec [N] 
	%\begin{itemize}
	%		\item $\widehat\varphi\circ\varphi = [\deg\varphi]_E$
	%		\item $\varphi\circ\widehat\varphi = [\deg\varphi]_{E'}$
	%	\end{itemize}
\end{frame}


\begin{frame}{$\insertframenumber -$SIDH problem}

	%\includegraphics[width=4.5cm]{media/FESTA-2.png}\\
	\begin{itemize}
		\item $d$ and $N$ smooth integers 
		\item secret isogeny $\varphi:E_0\to E_1$ of degree $d$
		\item P,Q a basis of the $N$-torsion of $E_0$
		\item $S = \varphi(P)$ and  $T=\varphi(Q)$
	\end{itemize}
	\begin{problem}
		Recover the secret isogeny $\varphi$ given curves $E_0,E_1$, points $P,Q,S,T$ and
		integers $d,N$.
	\end{problem}
	\vspace{0.5cm}
	\uncover<2>{\alert{Broken in polynomial time with a classic computer!}\\
	There exists a generic algorithm when $N^2>d$}

\end{frame}

\begin{frame}{$\insertframenumber -$Computationally isogeny with Scaled-Torsion Problem (CIST)}
	\begin{itemize}
		\item $d$ and $N$ smooth integers
		\item secret isogeny $\varphi:E_0\to E_1$ of degree $d$ 
		\item $P,Q$ a basis of the $N$-torsion of $E_0$
		\item diagonal matrix $\mathbf A$
			\[
			\begin{pmatrix} S \\ T \end{pmatrix} = \mathbf A \cdot\varphi
			\begin{pmatrix} P \\ Q \end{pmatrix} = 
			\begin{pmatrix} \lambda_1 & 0 \\ 0 & \lambda_2 \end{pmatrix} \cdot\varphi
			\begin{pmatrix} P \\ Q \end{pmatrix} 
			= \begin{pmatrix} [\lambda_1]\varphi(P) \\ [\lambda_2]\varphi(Q) \end{pmatrix} 
			\] 
	\end{itemize}
	\begin{problem}
		Recover the secret isogeny $\varphi$ given curves $E_0,E_1$, points $P,Q,S,T$, and
		integers $d,N$.
	\end{problem}
\end{frame}
%
%\begin{frame}{$\insertframenumber -$FESTA}
%	
%	Setup : curve $E_0$ and basis $P,Q$ of the $N$-torsion of $E_0$
%	\begin{itemize}
%		\item Alice computes secret keys $\varphi_A:E_0\to E_A$ and a matrix $\mathbf A$ \\
%		\item Alice computes the public key $(E_A,R_A,S_A)$ with
%			 \[
%			\begin{pmatrix} R_A \\ S_A \end{pmatrix} = \mathbf A\cdot\varphi_A
%			\begin{pmatrix} P \\ Q \end{pmatrix} 
%			\] 
%	\end{itemize}
%	\begin{itemize}
%		\item Bob computes secret keys $\varphi_1:E_0\to E_1$, $\varphi_2:E_A\to E_2$ and
%			matrix $\mathbf B$
%		\item Bob computes the public key $(E_1,R_1,S_1,E_2,R_2,S_2)$
%			\[
%		\begin{pmatrix} R_1 \\ S_1 \end{pmatrix} = \mathbf B\cdot\varphi_1
%		\begin{pmatrix} P \\ Q \end{pmatrix} \quad\quad
%		\begin{pmatrix} R_2 \\ S_2 \end{pmatrix} = \mathbf B\cdot\varphi_2
%		\begin{pmatrix} R_A \\ S_A \end{pmatrix} 
%			\]
%	\end{itemize}
%\end{frame}
\begin{frame}{$\insertframenumber -$FESTA}
	\hfill\raisebox{-40pt}[0pt][0pt]{\makebox[0pt][r]{
	%GRAPHIQUE
	\includegraphics[width=4.5cm]{media/FESTA-2.png}
}}
\begin{itemize}[<+->]
%\item Public parameters :\begin{itemize}
%		\item starting curve $E_0$
%		\item integers $d_A,d_1,d_2,N$
%		\item basis $P,Q$ of $E_0[N]$
%\end{itemize}
\item<1-> Trapdoor : $\varphi_A$ and matrix $\mathbf A$
\item<2-> Public key : $E_A$ and $R_A,S_A$  \[
			\begin{pmatrix} R_A \\ S_A \end{pmatrix} = \mathbf A\cdot\varphi_A
			\begin{pmatrix} P \\ Q \end{pmatrix} 
\] 



\item<3-> Input : $\varphi_1$,$\varphi_2$ and matrix $\mathbf B$
\item<4->	Output : curves $E_1,E_2$ and $R_1,S_1,R_2,S_2$  \[
		\begin{pmatrix} R_1 \\ S_1 \end{pmatrix} = \mathbf B\cdot\varphi_1
		\begin{pmatrix} P \\ Q \end{pmatrix}, \quad\quad
		\begin{pmatrix} R_2 \\ S_2 \end{pmatrix} = \mathbf B\cdot\varphi_2
		\begin{pmatrix} R_A \\ S_A \end{pmatrix} 
\] 

\end{itemize}
\begin{itemize}
	\item<2-> Alice : $\textsc{Trapdoor} \longrightarrow \textsc{Public Key}$
		\textcolor{ForestGreen}{(KeyGen)}
	\item<4-> Bob : $\textsc{Input} + \textsc{Public key} \longrightarrow \textsc{Output}$
		\textcolor{ForestGreen}{(Evaluation)}
	\item<5-> Alice : $\textsc{Output} + \textsc{Trapdoor} \longrightarrow \textsc{Input}$
		\textcolor{ForestGreen}{(Inversion)}
	\item<6-> $\textsc{Public key} \longrightarrow \textsc{Trapdoor}$
		\alert{(Hard)}
	\item<7-> $\textsc{Output}  \longrightarrow \textsc{Input}$
		\alert{(Hard)}
\end{itemize}
\end{frame}

\begin{frame}{$\insertframenumber -$FESTA Inversion (Alice)}
%	\hfill\raisebox{-40pt}[0pt][0pt]{\makebox[0pt][r]{
	%GRAPHIQUE
%}}
	\begin{center}	\includegraphics[width=4.5cm]{media/FESTA-2.png}\end{center}
	%TODO ajouter l'iso complète
	%TODO REFAIRE la slide
We have, if $\mathbf A$ and $\mathbf B $ commute,
	\begin{eqnarray*}
		[d_1]\begin{pmatrix} R_2 \\ S_2 \end{pmatrix} &=&
	\mathbf{B\cdot A\cdot B^{-1}}\cdot\varphi_2\circ\varphi_A\circ\widehat\varphi_1
	\begin{pmatrix} R_1 \\ S_1 \end{pmatrix} \\
																									&=& \mathbf A\cdot
																									\varphi_2\circ\varphi_A\circ\widehat\varphi_1
																									\begin{pmatrix} R_1 \\ S_1 \end{pmatrix}
\end{eqnarray*}
\begin{itemize}
	\item Matrix $\mathbf A$ is in $\textsc{Trapdoor}$
	\item We get an SIDH instance (no scaling) on isogeny
$\varphi_2\circ\varphi_A\circ\widehat\varphi_1$ 
	\item Using the attack against SIDH gives the isogeny
$\varphi_2\circ\varphi_A\circ\widehat\varphi_1$ 
\item Lead to recover $\textsc{Input}$ combined with the knowledge of $\varphi_A$
\end{itemize}

%	\begin{eqnarray*}
%		\psi \begin{pmatrix} R_1 \\ S_1 \end{pmatrix} 
%	& = & \psi \mathbf B\varphi_1 \begin{pmatrix} P \\ Q \end{pmatrix} 
%	= \varphi_2\circ\varphi_A [d_1]\mathbf B \begin{pmatrix} P \\ Q \end{pmatrix} \\
%	& = & %[d_1]\mathbf B\varphi_2\circ\mathbf{A^{-1}A}\varphi_A
%	\begin{pmatrix} P \\ Q \end{pmatrix} %\\
	 %= [d_1]\mathbf{BA^{-1}}\varphi_2\begin{pmatrix} R_A \\ S_A \end{pmatrix} 

%	\end{eqnarray*}
	%\includegraphics[width=4.5cm]{media/FESTA-2.png}
	
\end{frame}

\begin{frame}{$\insertframenumber -$Computationally isogeny with double Scaled-Torsion Problem (CIST2)}
	\begin{itemize}
		\item $d$,\textcolor{\mcolor}{$d'$} and $N$ smooth integers
		\item secret isogeny $\varphi:E_0\to E_1$ of degree $d$ 
		\item \textcolor{\mcolor}{secret isogeny $\varphi':E_0'\to E_1'$ of degree $d'$}
		\item $P,Q$ a basis of the $N$-torsion of $E_0$
		\item \textcolor{\mcolor}{$P',Q'$ a basis of the $N$-torsion of $E_0'$}
		\item diagonal matrix $\mathbf A$
			\[
			\begin{pmatrix} S \\ T \end{pmatrix} = \mathbf A \cdot\varphi
			\begin{pmatrix} P \\ Q \end{pmatrix}, 
			\quad
			\textcolor{\mcolor}{
			\begin{pmatrix} S' \\ T' \end{pmatrix} = \mathbf A \cdot\varphi'
			\begin{pmatrix} P \\ Q \end{pmatrix} 
		}
			\] 
	\end{itemize}
	\begin{problem}
		Recover the secret isogenies $\varphi$ and \textcolor{\mcolor}{$\varphi'$} given
		curves $E_0,E_1,\textcolor{\mcolor}{E_0',E_1'}$, points
		$P,Q,S,T,\textcolor{\mcolor}{P',Q'S',T'}$, and integers
		$d,\textcolor{\mcolor}{d'},N$.
	\end{problem}
\end{frame}

\section{Attack}

\begin{frame}{$\insertframenumber -$Problems}

	\begin{itemize}
		\item SIDH problem : Finding $\varphi:E_0\to E_1$ given  $E_0,E_1$ and
			$\varphi(P),\varphi(Q)$
			\textcolor{\mcolor}{too easy}
		\item CIST2 :Finding $\varphi,\varphi'$ given curves and scaled torsion point with the
			same matrix
		\item CIST : Finding $\varphi:E_0\to E_1$ fiven $E_0,E_1$ and
			$\lambda_1\varphi(P),\lambda_2\varphi(Q)$ 
		\item Isogeny problem : Finding $\varphi:E_0\to E_1$ given  $E_0,E_1$ \alert{too hard}
	\end{itemize}


\end{frame}
%TODO slide, CIST, CIST2 into FESTA
%\begin{frame}{$\insertframenumber -$Strategy}
%	\begin{center}\includegraphics[width=4.5cm]{media/FESTA-2.png}\end{center}
%	We want to solve the following problem :
%\begin{itemize}
%	\item $\textsc{Output} + \textsc{Public Key} \longrightarrow \textsc{Input}$ 
%\end{itemize}
%	If CIST is hard, it reduces to the two folling problems
%\begin{itemize}
%	\item $\textsc{Public Key} \longrightarrow \textsc{Trapdoor}$ : CIST \\
%	\item $\textsc{Output} \longrightarrow \textsc{Input}$ : CIST2 
%\end{itemize}	
%
%If an attack against CIST leaks relevant partial information for some instances, the
%previous reduction does not hold.
%
%%The hardness of CIST depends on starting curve $E_0$ and pubic parameters, there may be
%%situations where the public keys could leak relevant partial information on the trapdoor.
%
%\uncover<2>{\alert{$\rightarrow$ Motivates studying attacks against CIST}}
%
%%There exists recent work on the CIST problem \cite{polyCIST}
%\end{frame}
%\begin{frame}{$\insertframenumber -$Goals}
%	\begin{itemize}
%		\item 
%	Understand the lollipop attack and its upgraded version introduced by
%	Castryck and Vercauteren in \cite{polyCIST}. % TODO noms propres
%
%\item Get an effective algorithm that implements the attack with a clear statement on
%	correctness and running time.
%	\end{itemize}
%\end{frame}
\begin{frame}{$\insertframenumber -$Lollipop endomorphism}
	\hfill\raisebox{-100pt}[0pt][0pt]{\makebox[0pt][r]{
	%GRAPHIQUE
	\includegraphics[width=3cm]{media/lollipop-2.png}
}}
\begin{itemize}
	\item Inputs of the attack : $E,E'$
		\\$P,Q\in E[N]$
		$S,T\in E'[N]$
		\[
		\begin{pmatrix} S \\ T \end{pmatrix} =\mathbf A\cdot\varphi\begin{pmatrix} P \\ Q
	\end{pmatrix} 
	\quad\quad\quad\quad\quad
		\] 
\end{itemize}
\begin{enumerate}
		\item<2-> Find endomorphism $\omega$ and matrix $\mathbf M$
			\[
			\omega \begin{pmatrix} P \\ Q \end{pmatrix} = \mathbf M
			\begin{pmatrix} P \\ Q \end{pmatrix} 
			\] 
		\item<3-> Learn isogeny $\psi = \varphi\circ\omega\circ\widehat\varphi$	
			\[
				\psi\begin{pmatrix} S \\ T \end{pmatrix} = [\deg\varphi]\mathbf M
				\begin{pmatrix} S \\ T \end{pmatrix} 
			\] 
		\item<4-> Recover $\psi$ with SIDH attack \[
		N^2 > \deg\psi = \deg\varphi\deg\omega\deg\widehat\varphi = d^2\deg\omega
		\] 
		\item<5-> Recover $\varphi$ from $\psi$
	\end{enumerate}
\end{frame}

%TODO plus de slides la dessus…
\begin{frame}{$\insertframenumber -$Upgraded lollipop}
%TODO image
\end{frame}
\begin{frame}{$\insertframenumber -$Recovering $\varphi$ from $\psi$}
\end{frame}
\begin{frame}{$\insertframenumber -$Conclusion}
		
\end{frame}
\begin{frame}{$\insertframenumber -$References}
				\bibliographystyle{unsrt}
				%\bibliographystyle{amsalpha}
				\bibliography{biblio}
\end{frame}
\end{document}
